\section{Safety, Global, Societal, and Ethical Considerations}
\label{sec:impact}

\begin{table*}[!t]
\caption{Risks: Planned Mitigation and Outcome}
\label{tab:risks}
\centering
\footnotesize
\begin{tabular}{@{}p{0.15\textwidth}p{0.33\textwidth}p{0.45\textwidth}@{}}
\toprule
Risk & Planned mitigation (ECE 3992) & Outcome \\
\midrule
\multicolumn{3}{@{}l}{\emph{Mechanical}} \\
Hull integrity failure & Seven fiberglass plies, reinforced keel mount and stern, submersion test & Seven plies with reinforced keel boss and transom; leak test passed with no water inside \\
Keel detachment & Stainless bracket, bolts with lock nuts, embedded plate & 304 bracket bolted through the reinforced boss with nylon-insert lock nuts; zinc-plated fasteners corroding, to be replaced \\
Rudder bearing seizure & R188 ceramic bearings, marine grease, field-replaceable & Ceramic bearings in a removable two-part clamp \\
Mast or sail damage & Design margin; self-trimming wing feathers in gusts & Wing weathervanes freely; $\pm\SI{90}{\degree}$ stops protect the mast cable; high-wind check passed: the keel righted the boat after hard pushes and in field-test gusts \\
\midrule
\multicolumn{3}{@{}l}{\emph{Electrical}} \\
Water ingress & Sealed Pelican case, IP-rated pass-throughs, conformal coating & Printed enclosure under a latched hatch, cable glands, sealed SMA bulkhead and vane connector; breadboard not yet coated (open) \\
Battery failure & BMS & Per-cell voltage telemetry from four ADC taps; BMS fitted \\
LoRa link loss & \SI{10}{dBi} antennas, return to home, mission stored on board & Masthead \SI{10}{dBi} antenna; route stored on board with CRC; 95.9\% packet delivery to \SI{228}{\meter} in the field test; link-loss timeout implemented but off by default; return to home not implemented (open) \\
Servo failure & Waterproof 20 kg-cm servos, stall current monitoring & DS3218MG rudder servo; PID output limits and \SI{120}{\degree\per\second} slew limit; current monitoring not implemented \\
\midrule
\multicolumn{3}{@{}l}{\emph{Software}} \\
Sensor fusion divergence & Complementary filter, range checks, GPS-only fallback & BNO055 onboard fusion; every reading carries a validity code and staleness timeout; the heading PID stops when orientation data are unavailable \\
Real-time deadline miss & Task priorities, execution profiling, dual-core split & Radio moved to the CM4; fixed-rate tasks with consequence-ordered priorities \\
Path-planning edge cases & Tacking state machine, simulation before water trials & Tack planner implemented with 941 unit-test checks and a simulator route mode; not yet sailed on the water \\
No viable path & Planning timeout, loiter, then return to home & Corridor fault and arrival put the boat in SAFE; loiter and return to home not implemented (open) \\
\midrule
\multicolumn{3}{@{}l}{\emph{Found during development}} \\
Mismatched core firmware & -- & Shared-memory layout version checked at boot; CM4 refuses to attach on mismatch \\
Angle averaging at the seam & -- & Unwrapped running averages for encoder and vane; circular mean in navigation \\
\bottomrule
\end{tabular}
\end{table*}


\subsection{Public Health, Safety, and Welfare}
ALAN operates on open water near other watercraft, swimmers, and wildlife.
It is small (\SI{1}{\meter}, about \SI{5}{\kilo\gram}) and wind-powered, so
its speed and the energy of any collision are low. Testing is limited to
small, controlled bodies of water with an operator in sight, and the boat
reports its GPS position every second so it can always be located. The
operator can take control at any time with a mode command, and the SAFE mode
centers the rudder and feathers the tail. Within the firmware, each actuator
has a single writer with a timeout: if a controller crashes or is starved, the
rudder centers within \SI{200}{\milli\second} and the tail returns to its
safe angle within \SI{300}{\milli\second}. The link-loss timeout that drops
manual and semi-autonomous modes to SAFE is implemented but disabled for bench
work. It must be enabled before any unattended run. The 4S LiPo battery is the
main hazard on board; per-cell voltages are reported in every telemetry
packet so a weak or over-discharged cell is visible from shore.
The pack is fitted with a BMS and is charged with a professional battery
charger. Between uses it is unplugged completely from the circuit and left to
sit.

\subsection{Environmental Factors}
Wind propulsion produces no emissions and no underwater noise, and the battery
powers only the electronics. The hull is sealed fiberglass and epoxy, and the
keel and rudder hardware is stainless steel. Zinc-plated fasteners on the keel
are already corroding and will be replaced, so that no coating or rust
residue enters the water. A lost vessel would become debris, which is a
further reason for continuous position telemetry.

\subsection{Global and Cultural Factors}
Low-cost autonomous sailing platforms could make marine research more
accessible, particularly for coastal nations and institutions with limited
research budgets. The same technology could be misused, for example for
smuggling or for operation in protected marine areas without authorization.
Responsible deployment requires operating only where permitted.

\subsection{Economic Factors}
ALAN was built for roughly \$1000 from development boards, commodity sensors,
and fabrication resources available to most universities, a small fraction
of the cost of a commercial uncrewed sailing vessel. The firmware, simulator,
and design files are kept in a shared repository so other teams can build on
them.

\subsection{Professional and Regulatory Obligations}
Engineers have an obligation to ensure the vessel cannot harm people or the
environment, which motivates the failsafe behavior above, clear marking of
the boat as uncrewed, and operation only in permitted waterways. The radio
uses the unlicensed 902--928 MHz band, where Section 15.247 of the Federal Communications Commission (FCC) rules
covers only frequency-hopping and digitally modulated transmitters. In this
band, a hopping system with a channel narrower than \SI{250}{\kilo\hertz}
must use at least 50 frequencies and dwell no more than \SI{0.4}{\second} on
any one in \SI{20}{\second}, and a digitally modulated system must have a
bandwidth of at least \SI{500}{\kilo\hertz}~\cite{cfr15247}. ALAN currently
transmits on a single fixed \SI{125}{\kilo\hertz} channel for about
\SI{246}{\milli\second} every second, so it meets neither condition at its
present \SI{17}{dBm} output. Before further field use, the team should either
add frequency hopping across the band or reduce power to the limits for
non-hopping transmitters. Neither change has been made yet,
and the November 18 field test used the single-channel configuration.
Frequency hopping is the preferred option because it keeps the present range.
The SX1276 includes a hopping mode that changes channel on a programmable
period~\cite{sx1276}, and the ground-station radio would need to follow the
same channel sequence. This change is listed in Section~\ref{sec:future}.